\subsection{可度量化空间的仿紧性}

下述定理加强了第 1 号命题 2 的结果：

定理 4. 每个可度量化空间都是仿紧的。

该定理是下述四个引理的推论。

引理 3. 设 \(\Re = (U_{\alpha})_{\alpha \in A}\) 是可度量化空间 \(X\) 的一个开覆盖。则存在一个序列 \((\mathfrak{S}_n)\) ，它由 \(X\) 的开子集组成的局部有限族构成，使得 \(\mathfrak{S} = \bigcup \mathfrak{S}_n\) 是 \(X\) 的一个加细 \(\Re\) 的开覆盖。

设 \(d\) 是 \(\mathbf{X}\) 上与其拓扑相容的度量。对每个 \(\alpha \in \mathbf{A}\) 及每个整数 \(n\) ，令 \(\mathbf{F}_{n\alpha}\) 表示诸 \(x \in \mathbf{U}_{\alpha}\) 中满足 \(d(x, \mathbf{X} - \mathbf{U}_{\alpha}) \geqslant 2^{-n}\) 者所成的集。由于 \(\mathbf{X} - \mathbf{U}_{\alpha}\) 是闭集，我们有 \(\mathbf{U}_{\alpha} = \bigcup_{n} \mathbf{F}_{n\alpha}\) 。把集 \(\mathbf{A}\) 良序化；对每个 \(\alpha \in \mathbf{A}\) 及每个整数 \(n\) ，令 \(\mathbf{G}_{n\alpha}\) 为诸 \(x \in \mathbf{F}_{n\alpha}\) 中满足 \(x \notin \mathbf{F}_{n+1,\beta}\) （对一切 \(\beta < \alpha\) ）者所成的集，并令 \(\mathbf{V}_{n\alpha}\) 为诸 \(y \in \mathbf{X}\) 中满足 \(d(y, \mathbf{G}_{n\alpha}) - 2^{-n-3}\) 者所成的集。\(\mathbf{V}_{n\alpha}\) 显然是开集；另一方面，\(\mathbf{V}_{n\alpha} \subset \mathbf{U}_{\alpha}\) ，因为对每个 \(y \in \mathbf{V}_{n\alpha}\) 存在 \(x \in \mathbf{G}_{n\alpha}\) 使 \(d(x, y) \leqslant 2^{-n-1}\) ，而由 \(x \in \mathbf{F}_{n\alpha}\) ，我们有

\[
d (y, \mathrm{X} - \mathrm{U} _ {\alpha}) \geqslant d (x, \mathrm{X} - \mathrm{U} _ {\alpha}) - d (x, y) \geqslant 2 ^ {- n - 1},
\]

从而 \(y \in \mathbf{U}_{\alpha}\) 。再者，对每个 \(x \in \mathbf{X}\) ，设 \(\alpha\) 是 \(\mathbf{A}\) 中使 \(x \in \mathbf{U}_{\alpha}\) 的最小指标；则存在整数 \(n\) 使 \(x \in \mathrm{F}_{n\alpha}\) ，而由 \(\alpha\) 的定义推出 \(x \in \mathrm{G}_{n\alpha}\) ，从而 \(x \in \mathrm{V}_{n\alpha}\) 。这表明，若令 \(\mathfrak{S}_n = (\mathrm{V}_{n\alpha})_{\alpha \in \mathbf{A}}\) ，则 \(\mathfrak{S} = \bigcup_{n} \mathfrak{S}_n\) 便是 X 的一个开

覆盖，且加细 \(\Re\) ；于是剩下要证明的是，诸族 \(\mathfrak{S}_n\) 中的每一个都局部有限。为此我们先证明不等式 \(d(\mathrm{G}_{n\alpha}, \mathrm{G}_{n\beta}) \geqslant 2^{-n-1}\) 对 \(\alpha \neq \beta\) 成立。设 \(\beta < \alpha\) ；那么若 \(x \in \mathrm{G}_{n\alpha}\) 且 \(y \in \mathrm{F}_{n\beta}\) ，则由定义有 \(x \notin \mathrm{F}_{n+1, \beta}\) ，因而 \(d(x, \mathrm{X} - \mathrm{U}_{\beta}) < 2^{-n-1}\) ，而 \(d(y, \mathrm{X} - \mathrm{U}_{\beta}) \geqslant 2^{-n}\) ，于是 \(d(x, y) \geqslant 2^{-n-1}\) ；由于 \(\mathrm{G}_{n\beta} \subset \mathrm{F}_{n\beta}\) ，断言得证。由此，利用 \(\mathrm{V}_{n\alpha}\) 与 \(\mathrm{V}_{n\beta}\) 的定义立即推出 \(d(\mathrm{V}_{n\alpha}, \mathrm{V}_{n\beta}) \geqslant 2^{-n-2}\) 。由最后这个不等式推出：对每个 \(z \in \mathrm{X}\) ，以 \(z\) 为中心、\(2^{-n-3}\) 为半径的开球至多与族 \(\mathfrak{S}_n\) 中的一个集相交；因而 \(\mathfrak{S}_n\) 是局部有限族，证明完毕。

引理 4. 设 \((\mathfrak{S}_n)\) 是拓扑空间 \(\mathbf{X}\) 中开集的局部有限族所成的序列，使得

\[
\mathfrak {S} = \bigcup_ {n} \mathfrak {S} _ {n}
\]

是 X 的覆盖。则存在 X 的一个局部有限（但不一定开的）覆盖 B ，它加细 S 。

设 \(\mathbf{E}_n\) 是 \(\mathbf{X}\) 中的开集，即 \(\mathfrak{S}_n\) 中所有集的并；

令 \(U_{n}\) 表示 \(\bigcup_{k=1}^{n}E_{k}\) ，并令 \(A_{n}\) 表示 \(U_{n}-U_{n-1}\) （约定 \(U_{0}=\varnothing\) ）。考察形如 \(V\cap A_{n}\) 的子集所成的集 B ，其中 \(V\in S_{n}\) 而 n 为任意整数；我们来证明 B 满足引理的条件。对每个 \(x\in X\) 存在整数 n 使 \(x\in A_{n}\) ，因为诸 \(A_{n}\) 构成 X 的一个分划；于是 \(x\in E_{n}\) ，从而存在 \(V\in S_{n}\) 使 \(x\in V\) ；故 \(x\in V\cap A_{n}\) ，这就证明了 B 是 X 的覆盖。显然，此覆盖加细 S 。另一方面，对每个 \(x\in X\) 存在整数 n 使 \(x\in U_{n}\) ；由于 \(U_{n}\) 是开集而诸 \(S_{m}\) 是局部有限族，故对每个 m 存在 x 的邻域 \(W_{m}\) ，它含于 \(U_{n}\) 且只与 \(S_{m}\) 中有限多个集相交；因而 x 的邻域 \(W=\bigcap_{m=1}^{n}W_{m}\) 只与 B 中有限多个集相交，因为 \(W\cap A_{p}=\varnothing\) 对 p\textgreater n 成立。于是 B 局部有限，引理 4 得证。

引理 5. 设 X 是正则空间，使得对 X 的每个开覆盖 R ，存在 X 的一个（不一定开的）局部有限覆盖 A 加细 R 。则对 X 的每个开覆盖 R ，存在 X 的一个局部有限闭覆盖 F 加细 R 。

设 \(\mathfrak{R}\) 是 X 的任一开覆盖。对每个 \(x \in \mathbf{X}\) ，存在开集 \(\mathbf{U} \in \mathfrak{R}\) ，它包含 \(x\) ，从而（由于 X 正则）存在一个开邻域 \(\mathbf{V}_x\) （它包含 \(x\) ），使 \(\overline{\mathbf{V}}_x \subset \mathbf{U}\) 。族 \(\mathfrak{B}\) 由诸 \(\mathbf{V}_x\) 组成，它是 X 的一个开覆盖，故由假设存在一个局部有限覆盖 \(\mathfrak{B}\) 加细 \(\mathfrak{B}\) 。设 \(\mathfrak{F}\) 为 \(\mathfrak{B}\) 中诸集的闭包所成的族。由于覆盖 \(\mathfrak{B}'\) （由诸 \(\overline{\mathbf{V}}_x\) 组成）比 \(\mathfrak{R}\) 更细，而 \(\mathfrak{F}\) 又比 \(\mathfrak{B}'\) 更细，故 \(\mathfrak{F}\) 是 X 的一个加细 \(\mathfrak{R}\) 的闭覆盖。此外 \(\mathfrak{F}\) 还是局部有限的，因为若一个开集不与集 \(\mathbf{B} \in \mathfrak{B}\) 相交，则它也不与其闭包 \(\overline{\mathbf{B}}\) 相交。

引理 6. 设 X 是豪斯多夫空间，使得对 X 的任一开覆盖 R ，存在 X 的一个加细 R 的局部有限闭覆盖 F 。则 X 是仿紧的。

设 \(\mathfrak{R}\) 是 X 的任一开覆盖。我们要证明存在 X 的一个加细 \(\mathfrak{R}\) 的局部有限开覆盖。设 \(\mathfrak{A}\) 是 X 的一个加细 \(\mathfrak{R}\) 的局部有限覆盖（闭与否均可）；对每个 \(x \in \mathbf{X}\) ，设 \(W_x\) 是 \(x\) 的一个开邻域，它只与 \(\mathfrak{A}\) 中有限多个集相交。族 \(\mathfrak{B}\) 由诸集 \(W_x\) 组成，它是 X 的一个开覆盖。设 \(\mathfrak{F}\) 是 X 的一个加细 \(\mathfrak{B}\) 的局部有限闭覆盖。对每个 \(A \in \mathfrak{A}\) ，设 \(U_A\) 是 \(\mathfrak{R}\) 中包含 A 的一个集，并设 \(C_A\) 为诸集 \(F \in \mathfrak{F}\) 中满足 \(A \cap F = \emptyset\) 者的并。由于 \(\mathfrak{F}\) 局部有限，\(C_A\) 在 X 中闭（第 I 章，§ 1，第 6 号，命题 4），因而 \(A' = U_A \cap (X - C_A)\) 是开集。由于 \(A \cap C_A = \emptyset\) 且 \(A \subset U_A\) ，故 \(A \subset A'\) ，而族 \(\mathfrak{A}'\) （由诸集 \(A'\) 组成，其中 A 取遍 \(\mathfrak{A}\) ）是 X 的一个开覆盖；再者，由于 \(A' \subset U_A \in \mathfrak{R}\) ，\(\mathfrak{A}'\) 加细 \(\mathfrak{R}\) 。剩下要证明 \(\mathfrak{A}'\) 局部有限。对每个 \(x \in X\) 存在 \(x\) 的一个邻域 T ，它只与 \(\mathfrak{F}\) 中有限多个集、设为 \(F_1, \ldots, F_n\) 相交。由于每个 \(F_i\) 含于形如 \(W_{y_i}\) 的一个集中，按定义 \(F_i\) 只与 \(\mathfrak{A}\) 中有限多个集相交；设 \(A_{ij} (i \leqslant j \leqslant s_i)\) 就是这些集。若 A 是 \(\mathfrak{A}\) 中异于诸 \(A_{ij} (i \leqslant i \leqslant n, i \leqslant j \leqslant s_i)\) 的一个集，则由诸定义推出 \(A'\) 不与任何 \(F_i\) 相交，因而不与 \(T \subset \bigcup_{i=1}^{n} F_i\) 相交。这就完成了引理 6 与定理 4 的证明。

